\section{Spectral construction and normalization of the modal branch}
\label{ms:section}

This section supplies the global hypotheses used by
Proposition~\ref{hs:modal-transfer}.  Frequencies in this section are
complex vacuum wavenumbers.  An outgoing resonance means the
analytic continuation of the vacuum outgoing Green function from
$\operatorname{Im}k>0$; no decay at real spatial infinity is imposed
when $\operatorname{Im}k<0$.  The material is still real, positive,
scalar, nondispersive, and equal to vacuum outside a bounded set.
The normalization below is a reciprocal Green-function residue norm.
It is not a positive all-space energy integral.

\subsection{A compact analytic pencil and continuity under rounding}

Let $B$ be a ball containing all material changes, let
$X=L^2(B;\C^3)$, and extend its elements by zero when applying
whole-space operators.  Write $\chi=\varepsilon_r-1$.  With
$\Phi_k(x)=e^{ik|x|}/(4\pi|x|)$, put
\begin{equation}
 K(k)f=\left.(k^2I+\nabla\nabla)\Phi_k*f\right|_B,
 \qquad L_\chi(k)=I-K(k)\chi.
 \label{ms:pencil}
\end{equation}
Multiplication by $\chi$ is included in the last expression.
The static longitudinal Fourier projection, compressed to $B$, is
denoted by $P$.  It satisfies $P=P^*$ and $0\leq P\leq I$.
The distributional identity
\begin{equation}
 K(k)=-P+R(k),\qquad
 R(k)=k^2\Phi_k I+\nabla\nabla(\Phi_k-\Phi_0),\qquad R(0)=0
 \label{ms:compact-split}
\end{equation}
holds as an identity of bounded operators on $X$.

\begin{lemma}[PROVED: analytic Fredholm reduction with a fixed material bound]
\label{ms:fredholm}
For every measurable $0\leq\chi\leq C$ supported in $B$,
\begin{equation}
 L_\chi(k)=A_\chi F_\chi(k),\quad
 A_\chi=I+P\chi,\quad
 F_\chi(k)=I-A_\chi^{-1}R(k)\chi
 \label{ms:factorization}
\end{equation}
has $A_\chi$ invertible with $\|A_\chi^{-1}\|\leq1+C$.
The operator $R(k)$ is an entire Hilbert--Schmidt family, and
$F_\chi$ is an analytic identity-minus-compact pencil with
$F_\chi(0)=I$.  Its noninvertible values are isolated, of finite
algebraic multiplicity.  At $k\ne0$, its kernel is exactly the
restriction to $B$ of the finite local-energy Maxwell resonances
for this material, with exterior field given by \eqref{ms:pencil}.

If $0\leq\chi_j\leq C$ and $\chi_j\to\chi$ almost everywhere
in this same ball, then $F_{\chi_j}(k)\to F_\chi(k)$ in operator
norm, locally uniformly in $k$, as do all $k$ derivatives.
\end{lemma}

\begin{proof}
The Fourier symbol of $\nabla\nabla\Phi_0$ is
$-\xi\xi^T/|\xi|^2$, including its distributional local term.
Taylor expansion of $e^{ikr}$ shows that each entry of the
remaining kernel in \eqref{ms:compact-split} is $O(1/r)$ as
$r\downarrow0$, uniformly for $k$ in compact subsets of $\C$.
The constant term $ik/(4\pi)$ in $\Phi_k-\Phi_0$ has zero Hessian.
The squared kernel is locally integrable in three dimensions;
on bounded $B\times B$ it is integrable.  Its parameter power
series and derivatives have the same local majorants on compact
parameter sets.  This proves the entire Hilbert--Schmidt assertion.

Put $s=\sqrt\chi$.  The positive selfadjoint operator $I+sPs$
has inverse of norm at most one, and direct multiplication gives
\begin{equation}
 A_\chi^{-1}=I-Ps(I+sPs)^{-1}s.
 \label{ms:static-inverse}
\end{equation}
This proves the stated bound and factorization.  Analytic Fredholm
theory applies, with the invertible value at zero fixing its
nontrivial alternative.  This application needs no small contrast,
smooth boundary, or uniform feature size.

For completeness, the volume representation is an exact Maxwell
equation.  In distributions,
$(\operatorname{curl}\operatorname{curl}-k^2)K(k)f=k^2f$.
Consequently a kernel vector extended by $E=K(k)(\chi E)$
satisfies $\operatorname{curl}\operatorname{curl}E
=k^2(1+\chi)E$.  Its curl is locally square integrable, by
the first-derivative Newton-potential estimate; equivalently this
follows from local cutoff testing.  Set $H=(ik)^{-1}\operatorname{curl}E$
in units with vacuum impedance one.  This gives the two first-order
Maxwell equations and the outgoing continued representation.
Conversely, subtracting this volume field from any such resonance
leaves an entire vacuum Maxwell field regular at the origin and
outgoing outside a ball.  In each vector spherical harmonic its
regular radial function and outgoing radial function have nonzero
Wronskian for $k\ne0$, so all coefficients vanish.  Thus the
restriction is a kernel vector.  A zero restriction gives the zero
global field.

Almost-everywhere convergence and the common bound give strong
convergence of the multiplication operators.  The resolvent identity
and \eqref{ms:static-inverse} give strong convergence of $A_{\chi_j}^{-1}$
with a uniform bound.  Dominated convergence of the integral kernels
gives $R(k)\chi_j\to R(k)\chi$ in Hilbert--Schmidt norm.
A uniformly bounded strongly convergent operator sequence converges
in norm after composition on the right with a fixed compact operator:
approximate that compact operator by finite-rank operators.
Apply this observation to $A_{\chi_j}^{-1}R(k)\chi$ and add the
Hilbert--Schmidt difference.  Compactness in the parameter $k$ makes
the convergence locally uniform; Cauchy's formula gives the same
claim for derivatives.
\end{proof}

All these arguments also hold on a fixed closed symmetry sector
preserved by $P$, $R(k)$, and the material multiplication.  We use
the axisymmetric electric sector: $E$ has only meridional components
and $H$ only the azimuthal component in cylindrical coordinates.
Rotationally invariant scalar material preserves this sector.
It is closed in $L^2$ and orthogonal to the axisymmetric magnetic
sector, both for the Hermitian pairing and for the unconjugated
bilinear pairing.  At the coordinate axis these statements are
understood in the distributional and almost-everywhere sense.

\begin{lemma}[PROVED: persistence of a simple resonance]
\label{ms:continuation}
Suppose the restriction of $F_\chi$ to such a fixed sector has an
isolated characteristic value $k_*$ of algebraic multiplicity one.
An analytic $L^\infty$ material perturbation has a locally
holomorphic continued wavenumber $k_\eta$ and nonzero eigenfield
$E_\eta\in X$.  Their values at zero are $k_*$ and a prescribed
nonzero eigenfield.  If instead bounded nonnegative materials
$\chi_j$ converge almost everywhere to $\chi$, and preserve this
sector, they have unique nearby simple characteristic values
$k_j\to k_*$ and eigenfields $E_j\to E$ in $X$, after one fixed
nonzero linear functional fixes their scale.
\end{lemma}

\begin{proof}
Choose closed complements of the one-dimensional kernel and
cokernel of $F_\chi(k_*)$.  The block between those complements
is invertible, and remains so in a small operator-norm neighborhood.
Block elimination reduces the kernel equation to a single scalar
analytic function $f(k)$, multiplied by invertible analytic block
operators.  Algebraic multiplicity one is exactly
$f(k_*)=0$, $f'(k_*)\ne0$.
For an analytic material perturbation, $A_\chi^{-1}$ continues
analytically by a Neumann series near the base material, so every
block is jointly analytic.  The scalar analytic implicit-function
theorem supplies $k_\eta$; back substitution supplies $E_\eta$.
For almost-everywhere material convergence, Lemma~\ref{ms:fredholm}
gives uniform convergence of the blocks and scalar functions on a
small closed disk.  Rouch\'e's theorem gives exactly one zero,
counting multiplicity.  It tends to $k_*$ and is simple.  Back
substitution and a linear functional nonzero on $E$ give the stated
field convergence.  This reduction must use $F$ for the
almost-everywhere convergence claim: the leading multiplication
in $L$ need not converge in operator norm.  Multiplication by $A$
is analytic and invertible and does not alter a characteristic
value's multiplicity or its right kernel.
\end{proof}

\subsection{A simple sphere resonance without a genericity assumption}

Choose an algebraic real refractive index
$1<n<\sqrt{\varepsilon_+}$, leaving a positive margin for the
subsequent insertions.  Rational numbers already provide this
choice for every $\varepsilon_+>1$.  In a sphere of radius $a$ and
permittivity $n^2$, write $z=ka$ and
$\psi_l(z)=zj_l(z)$, $\xi_l(z)=zh_l^{(1)}(z)$.
Matching the tangential electric and magnetic fields of the regular
interior electric multipole and the outgoing exterior multipole
gives, for $l\geq1$, the determinant
\begin{equation}
 D_l(z)=n\psi_l(nz)\xi_l'(z)-\xi_l(z)\psi_l'(nz).
 \label{ms:mie-denominator}
\end{equation}
Primes mean derivatives with respect to the displayed argument.
For example the interior tangential electric component is
proportional to $\psi_l'(nz)/(nz)$, and its tangential magnetic
component to $n\psi_l(nz)/(nz)$; their exterior analogues are
$\xi_l'(z)/z$ and $\xi_l(z)/z$.  Eliminating their two amplitudes
gives \eqref{ms:mie-denominator}.

\begin{lemma}[PROVED: simplicity and separation of electric sphere poles]
\label{ms:sphere-simple}
For algebraic $n>1$, every nonzero zero of $D_l$ is simple, and
$D_l$ and $D_j$ have no common nonzero zero if $l\ne j$.
In the axisymmetric electric sector each such zero is therefore a
simple Maxwell characteristic value.  There are $l=1$ poles
\begin{equation}
 z_m=\frac{\pi(m+1/2)}{n}
       -\frac{i}{2n}\log\frac{n+1}{n-1}+O(m^{-1}),
 \qquad m\longrightarrow\infty.
 \label{ms:pole-asymptotic}
\end{equation}
Their real parts are positive, imaginary parts negative, and
$\operatorname{Re}z_m/(-2\operatorname{Im}z_m)\to\infty$.
The $l=1,m_{\rm az}=0$ interior electric field has a nonzero axial
component at the center and on a sufficiently short axial interval.
\end{lemma}

\begin{proof}
The exact finite spherical-Hankel expansions\cite{dlmfExplicit} are
\begin{align}
 \xi_l(z)&=(-i)^{l+1}e^{iz}p_l(z),&
 z h_l^{(2)}(z)&=i^{l+1}e^{-iz}\widetilde p_l(z),\\
 p_l(z)&=\sum_{s=0}^l\frac{i^s(l+s)!}{2^s s!(l-s)!z^s},&
 \widetilde p_l(z)&=\sum_{s=0}^l\frac{(-i)^s(l+s)!}{2^s s!(l-s)!z^s}.
 \label{ms:hankel-polynomials}
\end{align}
Using $\psi_l=(\xi_l+zh_l^{(2)})/2$, the determinant, up to a
nonzero factor independent of its zero condition, is
\begin{equation}
 C_{l,+}(z)e^{inz}+C_{l,-}(z)e^{-inz},
 \label{ms:exp-polynomial}
\end{equation}
where, with $p=p_l$, $\widetilde p=\widetilde p_l$,
\begin{align}
 C_{l,+}&=(-i)^{l+1}\{i(n-1)p(z)p(nz)
                  +np(nz)p'(z)-p(z)p'(nz)\},\\
 C_{l,-}&=i^{l+1}\{i(n+1)p(z)\widetilde p(nz)
             +n\widetilde p(nz)p'(z)-p(z)\widetilde p'(nz)\}.
 \label{ms:rational-coefficients}
\end{align}
These are rational functions with algebraic coefficients.  They
cannot vanish together at $z\ne0$: the two corresponding equations
would annihilate both independent interior Riccati--Hankel
solutions.  Their nonzero Wronskian would then force both
$\xi_l(z)$ and $\xi_l'(z)$ to vanish, contradicting uniqueness
for their second-order radial equation.  Thus at any zero of $D_l$
both coefficients are nonzero and the equation is equivalently
\begin{equation}
 e^{2inz}=R_l(z):=-C_{l,-}(z)/C_{l,+}(z).
 \label{ms:rational-pole}
\end{equation}
The expansion $p_l(z)=1+i l(l+1)/(2z)+O(z^{-2})$ gives
\begin{equation}
 R_l(z)=(-1)^l\frac{n+1}{n-1}
       \left(1-\frac{i l(l+1)}{nz}+O(z^{-2})\right).
 \label{ms:rational-asymptotic}
\end{equation}
For different parity the leading constants differ; for the same
parity and distinct $l$ the coefficients of $1/z$ differ.  Hence
$R_l-R_j$ is never identically zero when $l\ne j$.

The Lindemann--Weierstrass theorem implies that $e^\alpha$ is
transcendental for every nonzero algebraic complex
$\alpha$\cite{bezivinRobba1989}.  If $D_l$ and $D_j$ had a common
nonzero zero $z$, the nontrivial rational equation
$R_l(z)=R_j(z)$ would make $z$ algebraic.  Equation
\eqref{ms:rational-pole} would make $e^{2inz}$ algebraic, a
contradiction.  If $D_l$ had a multiple nonzero zero, differentiation
of \eqref{ms:rational-pole} would give
$R_l'(z)=2inR_l(z)$.  This too is a nontrivial rational equation:
$R_l'/R_l\to0$ at infinity, whereas $2in\ne0$.
It again makes $z$ algebraic and yields the same contradiction.
No division by a zero coefficient was used at a pole.

Separation of variables decomposes the axisymmetric electric
sector into the degrees $l=1,2,\ldots$, with one azimuthal mode
at each degree.  Regularity at the origin and the outgoing condition
leave one radial amplitude on each side of the sphere.  Locally in
$k\ne0$, solving the inhomogeneous radial equations and matching
them reduces each block, by analytic invertible elimination, to
its two-by-two tangential matching matrix.  Its determinant is
\eqref{ms:mie-denominator} times a nonzero analytic factor.
Thus a simple zero is a one-dimensional simple characteristic
value of that block.  Lemma~\ref{ms:fredholm} excludes accumulation
at a finite nonzero $k$; the other angular blocks are invertible
near the selected zero.  The proved separation across $l$ gives
algebraic multiplicity one in the whole selected sector.
The familiar degeneracy across the other azimuthal indices in the
full sphere is irrelevant to this invariant-sector statement.

For $l=1$ let $r_\infty=-(n+1)/(n-1)$ and choose the centers
$z_m^0$ equal to the two leading terms in
\eqref{ms:pole-asymptotic}.  They obey $e^{2inz_m^0}=r_\infty$.
On the circle $|z-z_m^0|=C/m$, the modulus of
$e^{2inz}-r_\infty$ is at least $cC/m$ for large $m$, with
$c>0$ independent of $m$.  The rational correction
$R_1(z)-r_\infty$ is at most $C_0/m$ there.  Taking $C$ large
enough, Rouch\'e's theorem gives exactly one zero in this disk.
This proves \eqref{ms:pole-asymptotic}.  Finally
$\psi_1(u)=u^2/3+O(u^4)$.  The regular electric dipole vector
spherical wave therefore tends to a nonzero constant axial vector
at the center, for a nonzero radial amplitude.  The amplitude
cannot vanish without making the whole matched eigenfield zero.
Analyticity in the homogeneous core preserves the nonzero component
on a short axial interval.
\end{proof}

\subsection{A genuine reciprocal mode norm}

The bilinear pairing in the remainder of this section has no
complex conjugation.  Reciprocity gives $K(k)^T=K(k)$ on $X$.
For a simple characteristic value in the selected sector, define
\begin{equation}
 d(E,k,\chi)=-\int_B \chi E\cdot K'(k)(\chi E)\,dx,
 \qquad \mathcal N(E)=-\frac{k}{2}d(E,k,\chi).
 \label{ms:residue-norm}
\end{equation}
The derivative kernel is a bounded compact operator; these are
finite integrals of locally square-integrable fields.  Their values
are independent of which containing ball is used.

\begin{lemma}[PROVED: nonzero residue norm and its stability]
\label{ms:normalization}
At a simple nonzero Maxwell pole in the selected sector,
$\mathcal N(E)\ne0$.  Scaling $E$ by $a\in\C$ scales
$\mathcal N$ by $a^2$.  The reciprocal electric Green operator for
$\operatorname{curl}\operatorname{curl}-k^2\varepsilon_r$,
restricted to this sector, has residue
\begin{equation}
 \operatorname*{Res}_{k=k_*}G(k;x,y)
       =-\frac{E(x)\otimes E(y)}{2k_*\mathcal N(E)}.
 \label{ms:green-residue}
\end{equation}
The norm is continuous along the almost-everywhere material and
eigenfield convergence of Lemma~\ref{ms:continuation}, and
holomorphic along its analytic branches.  A local nonzero analytic
choice of scale makes $\mathcal N=1$.
\end{lemma}

\begin{proof}
The functional $\ell_E(v)=\int_B\chi E\cdot v$ is nonzero:
if $\chi E=0$, the volume equation makes $E=0$.  It lies in the
left kernel of $L_\chi(k_*)$, since
$\ell_E L_\chi=\langle\chi E,\cdot\rangle-
\langle K(k_*)\chi E,\chi\,\cdot\rangle=0$.
It is a nonzero functional also on the chosen sector, because that
sector is closed under complex conjugation and multiplication by
the real rotationally invariant $\chi$.  Its derivative pairing
$\ell_E L_\chi'(k_*)E$ is exactly $d$.
The scalar Schur reduction in Lemma~\ref{ms:continuation} shows
that this pairing is nonzero at a simple characteristic value;
otherwise the scalar zero would have multiplicity at least two.
The residue of $L_\chi^{-1}$ is therefore
$E\otimes\ell_E/d$.

The vacuum electric resolvent is $G_0(k)=k^{-2}K(k)$, with its
natural restriction when the source is supported in $B$.
The full electric field for a compact source is
$L_\chi(k)^{-1}G_0(k)f$ in $B$.  At the pole, reciprocity and
$K(k_*)\chi E=E$ give
\begin{equation}
 \ell_E G_0(k_*)f=k_*^{-2}\int E(y)\cdot f(y)\,dy.
\end{equation}
The same identity for sources outside $B$ follows directly from
the global volume representation and the symmetric kernel.
Multiplying the two residue factors yields
$E\otimes E/(k_*^2d)$, which is
\eqref{ms:green-residue}.  This is a genuine pole normalization
of a quasinormal mode, consistent with reciprocal Green-residue
normalizations\cite{muljarov2016}; the sign reflects the displayed
choice of Maxwell operator.  For a degenerate full-space pole the
formula describes the selected symmetry block, not a claim that
the entire Green residue has rank one.

Under rounding, $\chi_j E_j\to\chi E$ in $L^2$: split the
difference into the uniformly bounded multiplier times $E_j-E$
and dominated convergence applied to $(\chi_j-\chi)E$.
Operator-norm continuity of $K'(k)$ proves convergence of $d$
and $\mathcal N$.  Holomorphy and quadratic scaling follow from
the formulas.  A nonzero analytic square root of $\mathcal N$
on a sufficiently small parameter disk normalizes the field.
\end{proof}

\begin{lemma}[PROVED: normalized resonance continuity away from rounded tips]
\label{ms:rounding}
Fix the finite axisymmetric conical construction and a simple
sector pole with $\mathcal N=1$.  Round all its vertices
axisymmetrically within balls of radius $\rho$, leaving the material
values and all other interfaces fixed.  For sufficiently small
$\rho$, there is a simple sector pole $k_\rho\to k_*$ and
eigenfield with $\mathcal N(E_\rho)=1$ converging to $E$ locally
in $L^2$.  On any fixed compact set away from the vertices,
including either side of a fixed smooth interface, the fields
converge uniformly one-sided.
\end{lemma}

\begin{proof}
The rounded coefficients converge almost everywhere with fixed
support and the same positive material bound.  Apply
Lemmas~\ref{ms:fredholm}--\ref{ms:normalization} and choose the
square-root scale converging to the sharp field's scale.
The volume representation supplies local $L^2$ convergence also
on any larger fixed ball, by the same static and compact estimates.
For local regularity write $U_\rho=E_\rho-E$.  Its equation is
\begin{equation}
 \operatorname{curl}\operatorname{curl}U_\rho
 =k_\rho^2\varepsilon_\rho U_\rho
 +(k_\rho^2\varepsilon_\rho-k_*^2\varepsilon)E.
\end{equation}
The right-hand side tends to zero in local $L^2$ and $U_\rho$
does likewise.  Cutoff testing gives
$\operatorname{curl}U_\rho\to0$ in smaller interior $L^2$.
The constant-coefficient curl/div estimate for
$\operatorname{curl}U_\rho$ then gives convergence in interior
$H^1$.  The fixed regularized curl lifting gives a difference
$W_\rho-W\to0$ in $H^2$, and the corresponding mean-fixed
potentials converge in $H^1$.
On a neighborhood disjoint from the changing tips the coefficients
are identical.  The scalar equation for the potential difference
is $\operatorname{div}(\varepsilon\nabla(v_\rho-v))
=-\operatorname{div}(\varepsilon(W_\rho-W))$.
The piecewise $H^3$ scalar transmission estimate across a fixed
smooth positive interface, used in
Lemma~\ref{hs:three-rounding}, gives convergence of the potentials
in smaller piecewise $H^3$ sets.  Hence the electric fields converge
piecewise in $H^2$, and uniformly by the three-dimensional Sobolev
embedding.  At ordinary interior points the same conclusion follows
from the homogeneous interior estimates.
\end{proof}

\subsection{Closed modal hotspot statement}

\begin{theorem}[PROVED: arbitrarily many peaks in one normalized finite-device mode]
\label{ms:modal-packing}
Fix $\varepsilon_+>1$, a real target wavenumber $k_0>0$, a
nonempty open patch $\mathcal P$ of a plane, a positive integer
$N$, and numbers $\ell,M,Q_0>0$.  There is a finite smooth
piecewise-constant dielectric with vacuum exterior and
$1\leq\varepsilon_r\leq\varepsilon_+$, and one nonzero outgoing
Maxwell quasinormal mode $(E,H,k_*)$, such that
\begin{equation}
 \operatorname{Re}k_*=k_0,\qquad
 \operatorname{Im}k_*<0,\qquad
 \frac{\operatorname{Re}k_*}{-2\operatorname{Im}k_*}>Q_0,
 \qquad \mathcal N(E)=1.
 \label{ms:modal-parameters}
\end{equation}
There are $N$ disjoint three-dimensional observation balls of
diameter less than $\ell$, centered in $\mathcal P$ and with
orthogonal projections contained in $\mathcal P$.  In each ball
the one-sided electric-field maximum exceeds $M$ and is more than
twice the maximum on a fixed collar of its boundary.  A point
attaining each maximum can be chosen in $\mathcal P$.
Thus each peak's half-amplitude connected region stays in its
small ball.  No uniform support-size bound, minimum feature size,
nondegeneracy of the maximum, or positive energy fraction per peak
is asserted.  The insertion contrast may depend on $N$; the common
material bounds do not.
\end{theorem}

\begin{proof}
Take an algebraic index $n$ strictly between one and
$\sqrt{\varepsilon_+}$.  Lemma~\ref{ms:sphere-simple} gives a
simple axisymmetric electric dipole pole with quality factor
strictly greater than $Q_0$.  Scale the sphere initially to make
the real part of its wavenumber $k_0$.  Put its symmetry axis in
the target plane and its center in the open patch.  Choose $N$
distinct points on a sufficiently short axial interval where the
background axial component is nonzero, with ample clearance from
the patch boundary.  Choose disjoint arbitrarily small conical
neighborhoods and close their material insertions smoothly away
from their vertices, inside the homogeneous sphere.

Axisymmetry preserves the electric sector.  Lemma~\ref{ms:continuation}
gives one holomorphic pole and eigenfield under simultaneous
positive cone contrast $\eta$.  Theorem~\ref{hs:modal-coefficient-continuation}
and Corollary~\ref{hs:modal-all-coefficients} give nonzero singular
coefficients at all $N$ vertices for one sufficiently small
$\eta>0$.  Choose it also so that $n^2(1+\eta)<\varepsilon_+$,
the pole remains simple, its imaginary part remains negative, and
its quality factor remains strictly above $Q_0$.
Lemma~\ref{ms:normalization} fixes its genuine norm to one.
Lemma~\ref{ms:rounding} supplies every rounding and normalization
hypothesis of Proposition~\ref{hs:modal-transfer}.  Its high-point
and collar argument now proves all the simultaneous strict peak
inequalities for sufficiently small rounding radius.  The same
axisymmetric rotation argument puts an actual maximizing point
of each ball in the specified plane.

There is one final harmless adjustment because the real pole
frequency has shifted.  About a point in the interior of the patch,
dilate space by $s=\operatorname{Re}k_*/k_0>0$.  Maxwell scaling
sends the wavenumber to $k_*/s$, keeps the quality factor and
material values unchanged, and tends to the identity as the
insertion contrast tends to zero.  A normalized field scales as
$E_s(x)=s^{-3/2}E(x/s)$: indeed \eqref{ms:residue-norm} and
the scaled kernel give $\mathcal N(E(x/s))=s^3\mathcal N(E)$.
Choose all observation-ball diameters with strict margins below
$\ell$, their projections with strict interior clearance, and
their pre-scaling peak targets larger than $s^{3/2}M$ for every
$s$ in a fixed small neighborhood of one.  These choices can be
made before the final rounding.  The dilation therefore preserves
every desired inequality and sets the exact real frequency.
\end{proof}

\begin{corollary}[PROVED: vanishing residue-normalized localization ratios]
\label{ms:modal-volume}
For one fixed sharp modal construction with $N$ tips, its rounded
normalized modes have uniformly finite nonzero residue norm, and
the maximum in each selected ball tends to infinity.  Consequently
\begin{equation}
 \frac{|\mathcal N(E_\rho)|}
 {\displaystyle\max_{x\text{ in ball }j}
        \varepsilon_\rho(x)|E_\rho(x)|^2}\longrightarrow0
 \quad (1\leq j\leq N).
\end{equation}
These are positive intensity ratios formed from the genuine
reciprocal norm.  A directional complex mode volume uses
$\mathcal N/[\varepsilon(u\cdot E)^2]$ instead; selecting a
Cartesian component of size at least $|E|/\sqrt3$ gives the same
vanishing magnitude along a subsequence with a fixed component.
No divergent all-space positive energy is substituted for
$\mathcal N$.
\end{corollary}

\begin{proof}
For every finite amplitude threshold choose ordinary high points
of the fixed sharp eigenfield at all tips, and then use the uniform
convergence at those fixed points from Lemma~\ref{ms:rounding}.
This proves divergence of every maximum as $\rho\downarrow0$.
The normalized numerator is one, and the positive permittivities
have a common lower bound.  There are only three Cartesian
components, so the indicated subsequence exists for a directional
volume.  The finite-frequency dilation has a finite nonzero limit
and leaves the vanishing-ratio conclusion unchanged.
\end{proof}
